\chapter{Polydipsia}

\section*{Definition}
Polydipsia is excessive thirst leading to abnormally high fluid intake, typically defined as consumption of \ensuremath{>}3 L of fluid per day in adults. It is often a presenting symptom of underlying metabolic or endocrine disorders and must be distinguished from primary (psychogenic) polydipsia.

\section*{What Happens in Your Body}
Polydipsia arises from osmotic or non-osmotic stimulation of the thirst center in the hypothalamic osmoreceptors. Hyperglycemia, hypernatremia, or hypercalcemia increases plasma osmolality, triggering antidiuretic hormone (ADH) release and thirst. In diabetes insipidus, impaired ADH secretion (central) or renal resistance (nephrogenic) leads to massive water loss and compensatory polydipsia. Primary polydipsia involves excessive voluntary fluid intake without osmotic stimulus.

\section*{Causes}
\begin{itemize}
  \item \textbf{Osmotic:} Diabetes mellitus (hyperglycemia), diabetes insipidus (central or nephrogenic)
  \item \textbf{Electrolyte disturbances:} Hypernatremia, hypercalcemia, hypokalemia
  \item \textbf{Renal:} Chronic kidney disease, diuretic use, post-obstructive diuresis
  \item \textbf{Psychogenic:} Schizophrenia, anxiety disorders, compulsive water drinking
  \item \textbf{Iatrogenic:} Anticholinergic medications, corticosteroids, lithium (nephrogenic DI)
\end{itemize}

\section*{When to See a Doctor}
See your doctor if:
\begin{itemize}
  \item Thirst is persistent, severe, or associated with polyuria (urine output \ensuremath{>}3 L/day)
  \item Accompanied by unintentional weight loss, blurred vision, or fatigue (suggests diabetes)
  \item Head trauma or neurosurgery preceding new-onset polydipsia
  \item Nocturnal thirst disrupting sleep
\end{itemize}

\section*{Related Symptoms}
\begin{itemize}
  \item Polyuria, nocturia, dry mouth, weight loss, fatigue, blurred vision, hypernatremia
\end{itemize}
\textbf{Important:} The water deprivation test remains the gold standard for distinguishing diabetes insipidus from primary polydipsia, though it must be performed under close supervision.

\section*{Self-Care / Home Management}
\begin{itemize}
  \item Track daily fluid intake and urine output to quantify the degree of polydipsia
  \item Avoid excessive caffeine and alcohol, which have diuretic effects
  \item Maintain a symptom diary to identify temporal patterns and associated symptoms
\end{itemize}

\section*{How Your Doctor Will Evaluate You}
\begin{itemize}
  \item Fasting plasma glucose, HbA1c, serum electrolytes (sodium, calcium, potassium), and renal function
  \item Urine osmolality and specific gravity on a random and first-morning sample
  \item Water deprivation test with serial serum and urine osmolality measurements; ADH levels as indicated
\end{itemize}

\section*{Treatment Options}
\begin{itemize}
  \item Glycemic control in diabetes mellitus with lifestyle modification, oral agents, or insulin
  \item Desmopressin (DDAVP) for central diabetes insipidus; thiazide diuretics and low-sodium diet for nephrogenic DI
  \item Correction of underlying electrolyte abnormalities and discontinuation of offending medications
  \item Behavioral therapy and fluid restriction for psychogenic polydipsia under close monitoring
\end{itemize}

\section*{References}
\begin{thebibliography}{99}
\bibitem{polydipsia:polydipsia1} American Diabetes Association. "Classification and diagnosis of diabetes: standards of medical care in diabetes-2022." Diabetes Care. 2022;45(Suppl 1):S17-S38.
\bibitem{polydipsia:polydipsia2} Robertson GL. "Diabetes insipidus." N Engl J Med. 2016;375(14):1387-1396.
\bibitem{polydipsia:polydipsia3} Verbalis JG. "Disorders of body water homeostasis." In: Harrison\'s Principles of Internal Medicine. 21st ed. McGraw-Hill; 2022.
\bibitem{polydipsia:polydipsia4} Schrier RW. "Body water homeostasis: clinical disorders of urinary concentration and dilution." J Am Soc Nephrol. 2006;17(7):1820-1832.
\bibitem{polydipsia:polydipsia5} Gardella A, Bhatt P. "Polydipsia and polyuria in children." Pediatr Rev. 2018;39(5):251-260.
\bibitem{polydipsia:polydipsia6} Berl T, Verbalis JG. "Pathophysiology of disorders of water metabolism." In: Brenner \& Rector\'s The Kidney. 11th ed. Elsevier; 2020.
\end{thebibliography}
